%! Inverse Functions — Unit 6, Lesson 6.7 — Warm-Up
\documentclass[10pt]{article}
\usepackage{algebra2-article}
\usepackage{algebra2-boxes}
\newcommand{\TallMath}[1]{$\displaystyle #1\rule[-1.4em]{0pt}{3.2em}$}

\begin{document}

\pageheader{Unit 6, Lesson 6.7}{Warm-Up}
\namedateperiod

\vspace{0.10in}

\begin{notesbox}{Getting Ready: Skills We'll Use Today}
\small Three short items. Every one of them is a piece of today's lesson --- item 2 is
\emph{yesterday's} discovery, and today it gets its name.

\vspace{4pt}
\textbf{1. Solve for the other variable (Units 1 and 6).} Each equation is solved for $y$. Solve it
for $x$ instead.
\begin{center}
\renewcommand{\arraystretch}{2.0}
\begin{tabularx}{0.96\linewidth}{@{}X X X@{}}
(a) \; $y=3x-7$ \quad $x=\blank{2.4cm}$ &
(b) \; $y=\dfrac{x+5}{2}$ \quad $x=\blank{2.4cm}$ &
(c) \; $y=x^{3}+1$ \quad $x=\blank{2.4cm}$ \\
\end{tabularx}
\end{center}
\vspace{-4pt}
Part (c) needed an operation from this unit. Which one, and why was it the right one to use?
\writeline

\vspace{0.12in}

\textbf{2. Both orders (Lesson 6.6).} Let $f(x)=2x+6$ and $g(x)=\dfrac{x-6}{2}$.
\begin{center}
\renewcommand{\arraystretch}{1.8}
\begin{tabularx}{0.96\linewidth}{@{}X X@{}}
(a) \; $(f\circ g)(x)=\blank{3.4cm}$ &
(b) \; $(g\circ f)(x)=\blank{3.4cm}$ \\
\end{tabularx}
\end{center}
\vspace{-6pt}
Yesterday you met one or two pairs that did this. In your own words: what are these two functions
doing to each other? \writeline

\vspace{0.12in}

\textbf{3. Undo the number trick.} A classmate says: \emph{``Think of a number. Double it. Then add
$5$. My answer is $17$.''}
\begin{center}
\renewcommand{\arraystretch}{1.6}
\begin{tabularx}{0.96\linewidth}{@{}X c c@{}}
\hline
 & \textbf{First move} & \textbf{Second move} \\
\hline
The trick did this, in this order & double & add $5$ \\
To get back, you do this, in this order & \blank{2.6cm} & \blank{2.6cm} \\
\hline
\end{tabularx}
\end{center}
\vspace{-2pt}
Their number was \blank{1.4cm}. \; To undo a two-step process you reverse each step \emph{and} you
also reverse the \blank{3.0cm}.

\end{notesbox}

\end{document}
